% 导言区
\documentclass{ctexart}

% \usepackage{ctex} % 中文处理宏包

% \begin{tabular}[<垂直对齐方式>][<列格式说明>]
    % <表项> & <表项> & <表项> & ... & <表项> \\
    % ......
% \end{tabular}
% 用\\表示换行
% 用&表示不同的列
% l-本列左对齐
% c-本列居中对齐
% r-本列右对齐
% p{<宽>}-本列宽度固定，能够自动换行


% 正文区
\begin{document}
    % tabular需要一个指定列排版格式的必选参数
    % 使用 l 左对齐
    % 使用 r 右对齐
    % 使用 c 居中对齐
    \begin{tabular}{l c c c r} % 本例中生成5个表格，分别是左对齐、居中对齐、居中对齐、居中对齐、右对齐
        姓名 & 语文 & 数学 & 外语 & 备注 \\
        张三 & 87 & 100 & 93 & 优秀 \\
        李四 & 75 & 64 & 52 & 补考 \\
        王二 & 80 & 82 & 78 & 良好 \\
    \end{tabular}

    \begin{tabular}{| l || c | c | c | r |}
        \hline
        姓名 & 语文 & 数学 & 外语 & 备注 \\
        \hline \hline
        张三 & 87 & 100 & 93 & 优秀 \\
        \hline
        李四 & 75 & 64 & 52 & 补考 \\
        \hline
        王二 & 80 & 82 & 78 & 良好 \\
        \hline
    \end{tabular}

    \begin{tabular}{| l || c | c | c | p{1.5cm} |} % 指定表格宽度，超过宽度自动换行
        \hline
        姓名 & 语文 & 数学 & 外语 & 备注 \\
        \hline \hline
        张三 & 87 & 100 & 93 & 优秀 \\
        \hline
        李四 & 75 & 64 & 52 & 补考另行通知 \\
        \hline
        王二 & 80 & 82 & 78 & 良好 \\
        \hline
    \end{tabular}
\end{document}
